\section{低支持度放宽肤色重选与分区形态学连接}\label{sec:low_support}

\subsection{内椭圆像素过少时的漏检现象}

在接通“双眼定位 + 内外双椭圆分区约束”的初版流程后，我们检查了 32 张验证图的结果。对于大多数在实验三中已经检出面部主体、只是外围粘连了头发或脖子的图像（例如 00611、00632），外椭圆裁剪与内椭圆闭运算已经能明显减少误检；但与此同时，有 4 张图像（00548、00562、00660、00733）的分割效果仍然很差，F1 停留在 $0.10\text{--}0.28$ 之间。

查看这 4 张图的中间掩膜可以发现，它们的共同问题是：在实验三 $k=1.25$ 阈值和最大连通域筛选之后，真正落入内椭圆 $E_{\text{in}}$ 的前景像素非常少（占内椭圆面积的比例 $\rho_{\text{in}}$ 都低于 $0.18$，例如 00562 在内椭圆里只剩下嘴角附近的一小块，00660 只选中了右下颌与脖子）。当内椭圆里几乎没有足够的初始肤色像素时，即使在内椭圆中做 $11\times11$ 闭运算和填孔，也没有东西可以连接，整张内脸依然是大片空白。

\subsection{内椭圆放宽肤色阈值重选与受限重建}

为什么这几张图在 $k=1.25$ 下落入内椭圆的像素这么少？回到第~\ref{sec:intro}~章提到的训练集肤色统计可以看到，$k=1.25$ 的矩形阈值本来只覆盖了约六到七成的训练肤色像素；当图像整体偏暗、受冷光影响或面部存在较大阴影时，真实脸部皮肤的 Cr、Cb 值很容易落到 $k=1.25$ 之外，但在 $k=2.00$（训练集肤色覆盖率达 91.86\%）下其实仍能被检出。

在实验三中，我们不敢直接把全图阈值放宽到 $k=2.00$，因为全图放宽会把大片头发和背景一起变成前景。但现在已经有了由双眼定出的内椭圆 $E_{\text{in}}$，情况就不同了：如果只在内椭圆像素明显不足的图像上，并且只允许在内椭圆 $E_{\text{in}}$ 内部使用放宽后的肤色阈值重选像素，外围背景和远端头发就不会被放进来。

按照这一思路，我们将“核心种子提取”和“低支持度放宽阈值重建”整理为以下三个步骤：
\begin{enumerate}
\item \textbf{提取高可信核心种子并计算内椭圆支持率：}设实验三输出的掩膜为 $M_{\text{Lab3}}$，先用外椭圆裁剪得到 $M_{\text{local}} = M_{\text{Lab3}} \cap E_{\text{out}}$，并在其中提取与内椭圆 $E_{\text{in}}$ 重叠像素数不少于 $N_{\text{core}}=24$ 的 8 连通分量，合记为 $M_{\text{sel}}$；若没有任何分量满足条件，则退一步保留 $M_{\text{local}}$ 中的最大连通块。将 $M_{\text{sel}}$ 落在内椭圆中的部分记为核心种子 $S = M_{\text{sel}} \cap E_{\text{in}}$，并计算内椭圆支持率
\begin{equation}
\rho_{\text{in}} = \frac{|S|}{|E_{\text{in}}|}.
\label{eq:innersupport}
\end{equation}
\item \textbf{判断低支持度并引入放宽肤色候选：}当 $\rho_{\text{in}} \ge 0.18$ 时，说明内椭圆中已有较多初始肤色像素，无需额外重选；当 $\rho_{\text{in}} < 0.18$ 时，判定该图内椭圆像素过少，计算阈值系数放宽至 $k_{\text{rec}}=2.00$ 时的肤色二值掩膜 $M_{\text{relax}}$，并截取其内椭圆部分 $M_{\text{relax}} \cap E_{\text{in}}$。
\item \textbf{闭运算连接与核心种子重叠过滤：}在 $E_{\text{in}}$ 内部，对 $M_{\text{relax}} \cap E_{\text{in}}$ 执行一次较大尺度的椭圆核闭运算（最终选用 $27\times27$ 核），把分散的脸部皮肤块连接成片，再与核心种子 $S$ 取并集：
\begin{equation}
A = \bigl(\operatorname{Close}_{27}(M_{\text{relax}}\cap E_{\text{in}})\cap E_{\text{in}}\bigr)\cup S.
\label{eq:allowmask}
\end{equation}
考虑到即使在内椭圆内部，放宽阈值后额头边缘也可能出现少量孤立的刘海误检块，我们对 $A$ 做 8 连通域分解，只保留与初始核心种子 $S$ 重叠像素数不少于 $24$ 的连通分量，最后在 $E_{\text{in}}$ 内部做一次漫水填孔，得到重建掩膜 $M_{\text{rec}}$：
\begin{equation}
M_{\text{rec}} = \operatorname{FillHoles}\!\left(\bigcup_{C_j\in\mathcal{C}_8(A):\,|C_j\cap S|\ge 24} C_j\right)\cap E_{\text{in}}.
\label{eq:lowsupport}
\end{equation}
\end{enumerate}

这样处理后，新增像素一方面被严格限制在内椭圆 $E_{\text{in}}$ 内部，另一方面必须通过放宽后的肤色闭运算与原有的高可信种子 $S$ 连成一体才会被保留，避免了不看图像肤色、直接把整个内椭圆涂满的做法。

\subsection{内外分区形态学操作与整体连接流程}

无论是常规样本还是触发了低支持重建的样本，最后都需要将内椭圆核心区与内外椭圆之间的过渡环带合成为完整的人脸掩膜。按第~\ref{sec:dual_ellipse}~章划定的三个空间区域，分区形态学连接的具体执行顺序如下：
\begin{enumerate}
\item \textbf{内椭圆核心区处理：}对保留的核心候选 $M_{\text{sel}}$ 执行 $11\times11$ 椭圆核闭运算并填孔，截取其落在内椭圆 $E_{\text{in}}$ 内的部分；若当前样本触发了低支持度重建（$\rho_{\text{in}}<0.18$），则将重建结果 $M_{\text{rec}}$ 一并并入核心区掩膜 $M_{\text{strong}}$：
\begin{equation}
M_{\text{strong}} = \bigl(\operatorname{FillHoles}(\operatorname{Close}_{11}(M_{\text{sel}}))\cap E_{\text{in}}\bigr) \cup M_{\text{rec}}.
\label{eq:mstrong}
\end{equation}
\item \textbf{过渡环带处理：}对 $M_{\text{sel}}$ 执行 $5\times5$ 椭圆核弱闭运算，仅截取其落在过渡环带 $E_{\text{out}}\setminus E_{\text{in}}$ 内的部分，记为 $M_{\text{weak}}$：
\begin{equation}
M_{\text{weak}} = \operatorname{Close}_{5}(M_{\text{sel}})\cap (E_{\text{out}}\setminus E_{\text{in}}).
\label{eq:mweak}
\end{equation}
\item \textbf{合并、种子连通过滤与最终填孔：}将核心区与过渡环带合并为候选掩膜 $M_{\text{prop}} = M_{\text{strong}} \cup M_{\text{weak}}$。为防止过渡环带中残留与面部主体断开的孤立碎块，再次对 $M_{\text{prop}}$ 提取 8 连通域，仅保留与核心种子 $S$ 相交的连通分量，最后做一次内部填孔并限制在 $E_{\text{out}}$ 内，得到最终分割掩膜 $M_{\text{final}}$：
\begin{equation}
M_{\text{final}} = \operatorname{FillHoles}\!\left(\bigcup_{C_j\in\mathcal{C}_8(M_{\text{prop}}):\,C_j\cap S\neq\emptyset} C_j\right)\cap E_{\text{out}}.
\label{eq:zonalmorph}
\end{equation}
\end{enumerate}

代码~\textbf{5-1} 列出了低支持度重建与内外分区形态学连接的核心实现。

\noindent\textbf{代码 5-1}\quad 低支持放宽阈值重建与双椭圆分区形态学连接（dual\_ellipse.py）
\begin{lstlisting}
def reconstruct_low_support(relaxed, seed, inner, kernel_size, minimum_seed_pixels):
    allowed = morphology(relaxed & inner, cv2.MORPH_CLOSE, kernel_size) & inner
    allowed |= seed & inner
    count, labels = cv2.connectedComponents(allowed, connectivity=8)
    retained = [c for c in range(1, count)
                if np.count_nonzero((labels == c) & seed.astype(bool)) >= minimum_seed_pixels]
    return fill_holes(np.isin(labels, retained).astype(np.uint8)) & inner

# 分区形态学连接：内椭圆强闭运算+低支持重建，过渡区弱闭运算，种子连通锚定
weak = morphology(selected, cv2.MORPH_CLOSE, config["weak_kernel"]) & outer
strong = fill_holes(morphology(selected, cv2.MORPH_CLOSE, config["strong_kernel"])) & inner
if relaxed is not None and inner_support < config["low_support_threshold"]:
    strong |= reconstruct_low_support(relaxed, seed, inner,
                                      config["reconstruction_close_kernel"],
                                      config["minimum_reconstruction_seed_pixels"])
proposal = strong | (weak & (outer & (1 - inner)))
count, labels = cv2.connectedComponents(proposal, connectivity=8)
anchor_labels = np.unique(labels[seed.astype(bool)])
final = fill_holes(np.isin(labels, anchor_labels[anchor_labels > 0]).astype(np.uint8)) & outer
\end{lstlisting}
